\documentclass[11pt,a4paper]{article}

%====================================================
% PACKAGES
%====================================================
\usepackage[utf8]{inputenc}
\usepackage[T1]{fontenc}
\usepackage[french]{babel}
\usepackage[a4paper,margin=1.7cm]{geometry}
\usepackage{amsmath,amssymb,siunitx}
\usepackage{lmodern}
\usepackage{enumitem}
\usepackage{fancyhdr}
\usepackage{lastpage}
\usepackage{hyperref}
\usepackage{xcolor}
\usepackage{array}
\usepackage{tabularx}
\usepackage{multicol}
\usepackage{graphicx}
\usepackage{tikz}
\usepackage{pgfplots}
\pgfplotsset{compat=1.18}
\sisetup{locale=FR}

%====================================================
% MISE EN PAGE
%====================================================
\setlength{\headheight}{14pt}
\pagestyle{fancy}
\fancyhf{}
\fancyhead[L]{Terminale générale -- Spécialité Physique-Chimie}
\fancyhead[C]{Devoir surveillé -- Mécanique}
\fancyhead[R]{Page \thepage/\pageref{LastPage}}
\fancyfoot[C]{\textit{Sujet original long avec barème détaillé}}

%====================================================
% COMMANDES
%====================================================
\newcommand{\vect}[1]{\overrightarrow{#1}}
\newcommand{\e}[1]{\vec{e}_{#1}}
\newcommand{\bareme}[1]{\hfill{\color{blue!70!black}\textbf{(#1 pts)}}}
\newcommand{\important}[1]{\textbf{#1}}

%====================================================
% DOCUMENT
%====================================================
\begin{document}

\begin{center}
    {\LARGE \textbf{Devoir surveillé de Physique-Chimie}}\\[0.2cm]
    {\large Terminale générale -- Spécialité}\\[0.2cm]
    \textbf{Chapitre : Mouvement et interactions -- Mécanique}\\[0.3cm]
\end{center}

\vspace{0.2cm}

\noindent\textbf{Consignes générales}
\begin{itemize}[leftmargin=0.7cm]
    \item La calculatrice est autorisée.
    \item Toute réponse doit être justifiée. Les résultats numériques doivent être donnés avec une unité cohérente.
    \item Le soin, la rigueur du raisonnement, la qualité des schémas, des bilans de forces et des projections seront valorisés.
    \item Le barème est détaillé question par question. Les questions demandant une vraie modélisation, une démarche longue ou une interprétation fine rapportent davantage que les questions d’application immédiate du cours.
    \item On pourra admettre un résultat démontré précédemment pour poursuivre.
\end{itemize}

\vspace{0.2cm}

\noindent\textbf{Données utiles}
\[
g=\SI{9.81}{m.s^{-2}}
\qquad
1~\text{km.h}^{-1}=\frac{1000}{3600}~\text{m.s}^{-1}
\]

\vspace{0.2cm}
\hrule
\vspace{0.3cm}

\begin{center}
\textbf{\large Exercice 1 -- Problème type bac métropole contextualisé}\\
\textbf{Sécurisation d’une piste de saut à ski nautique}
\end{center}

Une base de loisirs souhaite étudier la sécurité d’une rampe de saut utilisée pour le ski nautique. Un skieur quitte la rampe au point $O$ avec une vitesse initiale supposée horizontale. On étudie le mouvement du centre d’inertie $G$ du skieur après qu’il a quitté la rampe.

Le référentiel terrestre est supposé galiléen.

Le point $O$ est choisi comme origine des positions. On travaille dans un repère orthonormé $(O;\e{x},\e{z})$ où :
\begin{itemize}[leftmargin=0.7cm]
    \item l’axe $(Ox)$ est horizontal, orienté dans le sens du mouvement ;
    \item l’axe $(Oz)$ est vertical, orienté vers le haut.
\end{itemize}

À l’instant $t=0$, le skieur quitte la rampe depuis la hauteur
\[
h=\SI{2.40}{m}
\]
avec une vitesse de valeur
\[
v_0=\SI{54}{km.h^{-1}}.
\]

On néglige dans un premier temps les frottements de l’air.

\medskip

\textbf{Partie A -- Mise en équation du mouvement}

\begin{enumerate}[label=\textbf{A.\arabic*.}, leftmargin=0.9cm]
    \item Définir précisément le système étudié et le référentiel choisi. \bareme{1}
    \item Faire le bilan des forces exercées sur le skieur après son départ de la rampe. \bareme{2}
    \item En appliquant la deuxième loi de Newton, établir les coordonnées du vecteur accélération. \bareme{3}
    \item Déterminer les équations horaires $x(t)$ et $z(t)$ du mouvement. \bareme{4}
    \item Montrer que la trajectoire est parabolique et donner son équation cartésienne. \bareme{3}
\end{enumerate}

\medskip

\textbf{Partie B -- Exploitation numérique}

\begin{enumerate}[label=\textbf{B.\arabic*.}, leftmargin=0.9cm]
    \item Convertir la vitesse initiale en \si{m.s^{-1}}. \bareme{1}
    \item Déterminer la durée du vol jusqu’au contact avec l’eau, c’est-à-dire l’instant $t_f$ tel que $z(t_f)=0$. \bareme{3}
    \item En déduire la portée horizontale du saut. \bareme{2}
    \item Déterminer les coordonnées du vecteur vitesse à l’instant d’impact. \bareme{3}
    \item Calculer la valeur de la vitesse à l’impact. Comparer à la vitesse initiale et commenter physiquement. \bareme{2}
\end{enumerate}

\medskip

\textbf{Partie C -- Sécurité du bassin de réception}

Le bassin de réception commence à \SI{5.0}{m} horizontalement après l’extrémité de la rampe et s’étend sur une longueur de \SI{18.0}{m}.

\begin{enumerate}[label=\textbf{C.\arabic*.}, leftmargin=0.9cm]
    \item Le skieur tombe-t-il dans la zone prévue ? \bareme{2}
    \item On souhaite que le point d’impact soit situé entre \SI{8.0}{m} et \SI{15.0}{m} de la rampe. La vitesse initiale fixée à \SI{54}{km.h^{-1}} est-elle adaptée ? \bareme{2}
    \item Déterminer la valeur minimale de la vitesse initiale horizontale nécessaire pour atteindre \SI{8.0}{m}. \bareme{4}
\end{enumerate}

\medskip

\textbf{Partie D -- Prise en compte simplifiée de l’air}

Dans cette partie seulement, on modélise l’action de l’air par une force horizontale constante opposée au mouvement :
\[
\vect{F}_{\text{air}}=-f\,\e{x}
\qquad \text{avec} \qquad
f=\SI{60}{N}.
\]
La masse du skieur équipé est
\[
m=\SI{75}{kg}.
\]

\begin{enumerate}[label=\textbf{D.\arabic*.}, leftmargin=0.9cm]
    \item Donner les nouvelles coordonnées du vecteur accélération. \bareme{2}
    \item Établir la nouvelle expression de $x(t)$. \bareme{3}
    \item Sans refaire toute l’étude verticale, expliquer pourquoi le temps de chute reste inchangé. \bareme{1}
    \item Estimer alors la nouvelle portée horizontale. \bareme{2}
    \item Comment l’introduction de cette force modifie-t-elle la sécurité du saut ? \bareme{1.5}
\end{enumerate}

\vspace{0.4cm}
\hrule
\vspace{0.4cm}

\begin{center}
\textbf{\large Exercice 2 -- Exercice type bac centres étrangers, plus ambitieux}\\
\textbf{Virage d’une voiture sur route circulaire}
\end{center}

Une voiture de masse
\[
m=\SI{1.20e3}{kg}
\]
aborde un virage plan que l’on modélise par un arc de cercle de rayon
\[
R=\SI{45}{m}.
\]

On suppose que la voiture roule à vitesse constante sur la trajectoire circulaire. On étudie le mouvement dans le référentiel terrestre supposé galiléen.

On note $v$ la valeur de la vitesse de la voiture.

\medskip

\textbf{Partie A -- Analyse qualitative}

\begin{enumerate}[label=\textbf{A.\arabic*.}, leftmargin=0.9cm]
    \item Le mouvement de la voiture est-il uniforme ? Est-il accéléré ? Justifier soigneusement. \bareme{2}
    \item Représenter, sur un schéma, le vecteur vitesse et le vecteur accélération en un point du virage. \bareme{1.5}
    \item Donner l’expression de la valeur de l’accélération dans un mouvement circulaire uniforme. \bareme{1}
\end{enumerate}

\medskip

\textbf{Partie B -- Bilan des forces et condition dynamique}

On suppose dans un premier temps que la route est horizontale.

\begin{enumerate}[label=\textbf{B.\arabic*.}, leftmargin=0.9cm]
    \item Faire le bilan des forces extérieures exercées sur la voiture. \bareme{2}
    \item Expliquer pourquoi la résultante verticale des forces est nulle. \bareme{1}
    \item Identifier la force responsable de la courbure de la trajectoire. \bareme{1}
    \item Établir la relation entre la valeur $F_{\text{adh}}$ de la force latérale d’adhérence et la vitesse $v$. \bareme{3}
\end{enumerate}

\medskip

\textbf{Partie C -- Vitesse maximale sans glisser}

On admet que la valeur maximale de la force latérale d’adhérence vaut
\[
F_{\max}=\mu m g
\]
avec
\[
\mu=0.65.
\]

\begin{enumerate}[label=\textbf{C.\arabic*.}, leftmargin=0.9cm]
    \item Déterminer la vitesse maximale $v_{\max}$ permettant à la voiture de suivre le virage sans glisser. \bareme{4}
    \item Exprimer le résultat en \si{km.h^{-1}}. \bareme{1}
    \item Une voiture entre dans ce virage à \SI{80}{km.h^{-1}}. Le virage est-il théoriquement franchissable dans ces conditions ? \bareme{2}
\end{enumerate}

\medskip

\textbf{Partie D -- Influence de la pluie}

Sous la pluie, le coefficient devient
\[
\mu'=0.35.
\]

\begin{enumerate}[label=\textbf{D.\arabic*.}, leftmargin=0.9cm]
    \item Déterminer la nouvelle vitesse maximale. \bareme{2}
    \item Comparer les deux vitesses maximales et commenter le risque routier. \bareme{1.5}
    \item Montrer que le rapport des vitesses maximales vaut
    \[
    \frac{v'_{\max}}{v_{\max}}=\sqrt{\frac{\mu'}{\mu}}.
    \]
    \bareme{2}
\end{enumerate}

\medskip

\textbf{Partie E -- Question plus fine de modélisation}

Le conducteur affirme : \emph{« Si je garde la même valeur de vitesse, je n’accélère pas. »}

\begin{enumerate}[label=\textbf{E.\arabic*.}, leftmargin=0.9cm]
    \item Expliquer précisément pourquoi cette phrase est fausse du point de vue de la mécanique. \bareme{2}
    \item Indiquer ce que mesure réellement l’accélération dans cette situation. \bareme{2}
\end{enumerate}

\vspace{0.4cm}
\hrule
\vspace{0.4cm}

\begin{center}
\textbf{\large Exercice 3 -- Problème plus difficile et plus théorique}\\
\textbf{Comment viser le plus loin possible ?}
\end{center}

On étudie le tir d’une balle dans un champ de pesanteur uniforme, sans frottements. La balle est lancée depuis le point $O$ situé au niveau du sol, avec une vitesse initiale de valeur $v_0$ faisant un angle $\alpha$ avec l’horizontale.

On travaille dans le repère orthonormé $(O;\e{x},\e{z})$, avec $\e{x}$ horizontal et $\e{z}$ vertical vers le haut.

Dans tout l’exercice,
\[
v_0=\SI{18.0}{m.s^{-1}}
\qquad \text{et} \qquad
g=\SI{9.81}{m.s^{-2}}.
\]

\medskip

\textbf{Partie A -- Équations horaires}

\begin{enumerate}[label=\textbf{A.\arabic*.}, leftmargin=0.9cm]
    \item Donner les coordonnées initiales du vecteur vitesse en fonction de $v_0$ et $\alpha$. \bareme{1}
    \item Établir les équations horaires :
    \[
    x(t)=v_0\cos\alpha\, t
    \qquad \text{et} \qquad
    z(t)=v_0\sin\alpha\, t-\frac12gt^2.
    \]
    \bareme{3}
    \item Retrouver l’équation cartésienne de la trajectoire. \bareme{3}
\end{enumerate}

\medskip

\textbf{Partie B -- Temps de vol et portée}

\begin{enumerate}[label=\textbf{B.\arabic*.}, leftmargin=0.9cm]
    \item Déterminer l’instant non nul auquel la balle revient au sol. \bareme{3}
    \item En déduire la portée horizontale :
    \[
    D(\alpha)=\frac{v_0^2}{g}\sin(2\alpha).
    \]
    \bareme{4}
    \item Vérifier l’homogénéité de cette expression. \bareme{1}
\end{enumerate}

\medskip

\textbf{Partie C -- Quel angle donne la portée maximale ?}

\begin{enumerate}[label=\textbf{C.\arabic*.}, leftmargin=0.9cm]
    \item À quelle condition sur $\sin(2\alpha)$ la portée est-elle maximale ? \bareme{1}
    \item En déduire l’angle $\alpha_{\max}$ donnant la portée maximale. \bareme{2}
    \item Déterminer la portée maximale correspondante. \bareme{2}
\end{enumerate}

\medskip

\textbf{Partie D -- Un résultat rigolo et symétrique}

On considère deux angles de tir $\alpha$ et $\beta$ tels que
\[
\alpha+\beta=90^\circ.
\]

\begin{enumerate}[label=\textbf{D.\arabic*.}, leftmargin=0.9cm]
    \item Montrer que les tirs d’angles $\alpha$ et $\beta$ ont la même portée. \bareme{2}
    \item Interpréter physiquement ce résultat. \bareme{1.5}
    \item Calculer les deux angles différents qui permettent d’atteindre une cible située à \SI{20.0}{m}. \bareme{4}
\end{enumerate}

\medskip

\textbf{Partie E -- Une vraie réflexion physique}

On cherche maintenant non plus à aller le plus loin, mais à atteindre une cible placée à l’abscisse
\[
x_c=\SI{12.0}{m}
\]
et à la hauteur
\[
z_c=\SI{3.00}{m}.
\]

\begin{enumerate}[label=\textbf{E.\arabic*.}, leftmargin=0.9cm]
    \item Écrire l’équation que doit vérifier $\alpha$ pour que la balle passe par le point $(x_c,z_c)$. \bareme{3}
    \item Expliquer pourquoi il peut y avoir, suivant les données, zéro, une ou deux trajectoires possibles. \bareme{2.5}
    \item Sans résoudre complètement l’équation, interpréter les deux solutions éventuelles du point de vue de la physique du tir. \bareme{2}
\end{enumerate}

\vspace{0.4cm}
\hrule
\vspace{0.4cm}

\begin{center}
\textbf{\large Exercice 4 -- Découverte d’une nouvelle notion}\\
\textbf{La force de recentrage dynamique}
\end{center}

Dans tout cet exercice, on introduit un modèle fictif, inspiré des forces de rappel, appelé \important{force de recentrage dynamique}. Cette force est définie, pour un mobile se déplaçant sur un axe horizontal $(Ox)$, par :
\[
\vect{F}_r=-k\,x\,\e{x}
\]
où :
\begin{itemize}[leftmargin=0.7cm]
    \item $x$ est l’abscisse du mobile ;
    \item $k$ est une constante strictement positive, appelée \important{raideur de recentrage}.
\end{itemize}

Cette force n’est pas au programme comme notion officielle : on l’introduit ici pour travailler la modélisation, le sens d’une force et l’interprétation des équations de Newton.

On considère un mobile de masse $m$ se déplaçant sur un rail horizontal sans frottement.

\medskip

\textbf{Partie A -- Compréhension de la nouvelle force}

\begin{enumerate}[label=\textbf{A.\arabic*.}, leftmargin=0.9cm]
    \item Donner la dimension physique du coefficient $k$ à partir de la relation
    \[
    F=kx.
    \]
    \bareme{1.5}
    \item Quel est le sens de la force si $x>0$ ? si $x<0$ ? \bareme{2}
    \item Que vaut la force lorsque le mobile passe par l’origine ? \bareme{1}
    \item Expliquer pourquoi le nom \emph{force de recentrage} est pertinent. \bareme{1.5}
\end{enumerate}

\medskip

\textbf{Partie B -- Conséquence dynamique}

\begin{enumerate}[label=\textbf{B.\arabic*.}, leftmargin=0.9cm]
    \item Appliquer la deuxième loi de Newton sur l’axe $(Ox)$ et établir une relation entre $a_x$ et $x$. \bareme{3}
    \item Si $x>0$, quel est le signe de l’accélération ? Interpréter physiquement. \bareme{2}
    \item Si $x<0$, quel est le signe de l’accélération ? Interpréter physiquement. \bareme{2}
    \item Montrer que l’accélération est toujours dirigée vers l’origine. \bareme{2}
\end{enumerate}

\medskip

\textbf{Partie C -- Étude qualitative du mouvement}

On écarte le mobile vers la droite puis on le lâche sans vitesse initiale.

\begin{enumerate}[label=\textbf{C.\arabic*.}, leftmargin=0.9cm]
    \item Décrire qualitativement l’évolution du mouvement juste après le lâcher. \bareme{2}
    \item Que se passe-t-il lorsque le mobile arrive à l’origine ? \bareme{2}
    \item Pourquoi peut-on s’attendre, qualitativement, à un mouvement d’aller-retour autour de l’origine ? \bareme{2.5}
\end{enumerate}

\medskip

\textbf{Partie D -- Étude quantitative sur quelques situations}

On prend dans cette partie :
\[
m=\SI{0.400}{kg}
\qquad \text{et} \qquad
k=\SI{5.0}{N.m^{-1}}.
\]

\begin{enumerate}[label=\textbf{D.\arabic*.}, leftmargin=0.9cm]
    \item Calculer la valeur de la force lorsque $x=\SI{0.10}{m}$. \bareme{1}
    \item Calculer la valeur de l’accélération correspondante. \bareme{2}
    \item Même question pour $x=\SI{0.25}{m}$. \bareme{2}
    \item Comparer les deux accélérations et commenter. \bareme{1.5}
\end{enumerate}

\medskip

\textbf{Partie E -- Comparaison avec le cours}

\begin{enumerate}[label=\textbf{E.\arabic*.}, leftmargin=0.9cm]
    \item Expliquer en quoi cette force diffère du poids. \bareme{1.5}
    \item Expliquer en quoi elle diffère d’une force de frottement. \bareme{1.5}
    \item Citer une situation réelle de physique où une force de rappel vers une position d’équilibre existe effectivement. \bareme{1.5}
\end{enumerate}

\medskip

\textbf{Partie F -- Réflexion plus fine}

On ajoute maintenant une force de frottement modélisée par :
\[
\vect{f}=-\lambda \vect{v}
\]
avec $\lambda>0$.

\begin{enumerate}[label=\textbf{F.\arabic*.}, leftmargin=0.9cm]
    \item Quel est l’effet qualitatif de cette force sur le mouvement d’aller-retour décrit plus haut ? \bareme{2}
    \item Peut-on s’attendre à ce que le mobile finisse par s’arrêter ? Si oui, où ? \bareme{2}
    \item Pourquoi cette version du modèle est-elle plus réaliste qu’un modèle sans frottements ? \bareme{1.5}
\end{enumerate}

\vspace{0.4cm}
\hrule
\vspace{0.4cm}

\begin{center}
\textbf{\large Barème global indicatif}
\end{center}

\begin{center}
\renewcommand{\arraystretch}{1.25}
\begin{tabular}{|>{\raggedright\arraybackslash}p{9.8cm}|c|}
\hline
\textbf{Compétence évaluée} & \textbf{Valorisation} \\
\hline
Identifier correctement système, référentiel, forces usuelles, donner une formule du cours & Faible à modérée \\
\hline
Projeter correctement la deuxième loi de Newton, établir une équation horaire, exploiter un modèle simple & Modérée \\
\hline
Mener une étude complète avec modélisation, interprétation, cohérence physique et calculs longs & Forte \\
\hline
Justifier finement un résultat qualitatif, discuter un modèle, commenter les limites d’une hypothèse & Forte \\
\hline
Réussir les questions ouvertes de réflexion physique et d’analyse critique & Très valorisé \\
\hline
\end{tabular}
\end{center}

\vfill

\begin{center}
\textbf{Fin du sujet}
\end{center}

\end{document}